% =============================================================================
% chapter5_discussion.tex — Chapter 5: Discussion
% =============================================================================
\chapter{Discussion}
\label{chap:discussion}

% TODO: Write full prose for each section.
% The scaffolding below provides the argumentative skeleton derived from
% the five experiments. Each section lists the key claims to make and
% the prior-work references to cite.

% ---------------------------------------------------------------------------
\section{Interpretation of Results}
\label{sec:discussion_interpretation}
% KEY ARGUMENT:
%   The five experiments collectively demonstrate that PQ optimisation via
%   L-BFGS-B with analytical gradients is PRACTICAL for architectural design
%   for meshes up to ~5,600 faces.
%
%   Connect each experiment back to the central research question:
%   "Can PQ mesh optimisation be made practical for architectural design?"
%
%   EXP-01: YES — interactive (<2 s) for up to 400 faces; automated batch
%            for up to 5,625 faces in ~80 s.
%   EXP-02: YES — hardware acceleration provides 2–3× additional speedup
%            at zero cost to numerical accuracy.
%   EXP-03: YES — convergence in 9–13 iterations is mesh-size-independent,
%            which is a critical property for predictable design workflows.
%   EXP-04: YES — the tool is robust to weight mis-specification; users
%            do not need expert calibration to obtain good results.
%   EXP-05: YES — 4/4 convergence on heterogeneous real-world meshes;
%            oloid result validates theoretical developable surface behaviour.

% TODO: Write §5.1 prose here.

% ---------------------------------------------------------------------------
\section{Comparison with Prior Work}
\label{sec:discussion_prior_work}
% KEY ARGUMENT:
%   The measured empirical complexity O(n^{1.27}) is competitive with
%   state-of-the-art methods but does not yet exploit sparsity structure.
%
%   COMPARE AGAINST:
%   - Liu et al. (2006): O(n log n) ARAP-based conical mesh optimisation;
%     exploits the sparsity of the as-rigid-as-possible energy Hessian.
%   - Pottmann et al. (2007): conical mesh methods via conjugate-net
%     parameterisation; geometric rather than energy-minimisation approach.
%   - Note: this tool uses dense batched SVD (no sparsity exploitation),
%     hence the n^{1.27} exponent instead of n log n.
%
%   CITE: Liu2006, Pottmann2007, Nocedal2006

% TODO: Write §5.2 prose here.

% ---------------------------------------------------------------------------
\section{Limitations}
\label{sec:discussion_limitations}
% FOUR FORMALLY DOCUMENTED LIMITATIONS:
%
%   L1 (from architecture.md + EXP-05):
%       High-valence vertex exclusion from angle-balance energy.
%       Vertices with valence > threshold are silently excluded from E_a,
%       which may produce artefacts on irregular quad meshes with
%       extraordinary vertices (e.g. the Bob mesh).
%
%   L2 (from EXP-05, oloid anomaly documentation):
%       Metric degeneracy for near-planar inputs.
%       The planarity improvement percentage (before-after)/before becomes
%       numerically meaningless when the initial planarity error is near
%       zero (oloid case). Absolute deviation must be used instead.
%
%   L3 (from methodology.md §4.3 + future_works.md):
%       No sparse L-BFGS variant.
%       The current L-BFGS-B implementation uses a dense Hessian
%       approximation. A sparse L-BFGS variant (Byrd et al., 1994) would
%       reduce complexity towards O(n log n) for structured meshes.
%
%   L4 (from methodology.md §5):
%       Degenerate-face area threshold is heuristic.
%       The preprocessor removes faces with area < 10^{-10} (unit-scale),
%       but this threshold is not formally justified. Meshes at non-unit
%       scale may require adjustment, and the threshold is not exposed
%       as a user-configurable parameter.
%
%   CITE: Byrd1994, Nocedal2006

% TODO: Write §5.3 prose here.
